% !TeX spellcheck = ru_RU-Russian
% Класс документа для стандартного формата
\documentclass[a4paper,12pt]{article}

% Поддержка русского языка
\usepackage[utf8]{inputenc}
\usepackage[russian]{babel}

% Поддержка математических формул
\usepackage{amsmath, amssymb}

% Работа с изображениями (для PNG)
\usepackage{graphicx}
\usepackage{tikz}
\usetikzlibrary{shapes.geometric, arrows, positioning, fit, calc}
\graphicspath{{./images/}} % Путь к папке с изображениями

% Форматирование страницы
\usepackage{geometry}
\geometry{margin=2cm} % Установка полей

% Поддержка списков
\usepackage{enumitem}

% Шрифты (Computer Modern, совместим с PDFLaTeX)
\usepackage[T1]{fontenc} % Кодировка шрифтов для латинских и кириллических символов

\begin{document}
\title{Определение оптимального решения задачи движения, распределения звёзд в диске Галактики и вычисление объективно значимых данных из этих моделей.}
\maketitle
\section{Работа с потенциалами, силами, скоростями, моделями минимизации.}
Полагаясь на две модели, что можно выделить из имеющихся у нас представлений, можно вывести 2 модели оптимизации на поиск изменения диска звёзд, являющейся единственной объективно значимой задачей оптимизации на данный момент. Первое представление - цефеиды непосредственно и образуют этот диск. Тогда поведение звёзд будет определяться динамически - сумма сил будет определяться из стоячих потенциалов + изменяющихся во времени и пространстве (не только по двум осям R и Z). Второе предположение основано на распределении энергий звёзд в пространстве:\\
\textit{$\bullet$ Во-первых, распределение энергий в системе несталкивающихся тел определяется случайным распределением для газов (к примеру), чем мы и воспользуемся}\\
\textit{$\bullet$ Во-вторых, распределение энергий также задаётся видом дисков, но эта энергия - потенциальная - учтём её с кинетической, чтобы найти полную энергию}

Этот подход может быть полезен в рассмотрении факта о том, что звёзды разлетаются в разных направлениях при столкновении, но сумма их энергий продолжает иметь константный вид. Это влияние можно будет учесть, к примеру, наложением случайного распределения, вложенного в случайное распределение по энергиям.

На этом представления могут быть исчерпаны или же сами результаты, получаемые из данных новых положений, могут не нести объективную значимость.

\subsection{Модель потенциала.}
Искомое распределение будет рассчитываться для скорости в особой конфигурации системы - в состоянии, когда диск и остальные потенциалы ориентированы вдоль плоскости XY. При этом расстояние для них берётся с условием (к примеру, ЗСМИ) (радиус вектор строго определён для определения предположительного начального расстояния звезды).

Для такой модели будет определён ЗСЭ:

\begin{equation}
	V_{now}^2/2 + U(X, Y, Z) = V_{x}^2/2 + U(R_x, Z_0 = 0),
\end{equation}
где $R_x$ может быть приближён уравнением $\sqrt{X^2 + Y^2}$ для звёздного диска в плоскости XY.


\section{Как работает программа.}
\subsection{Поиск перпендекуляров к кривой полиномиальной поверхности.}

\textit{Введём 4 точки относительно начальной заданной ref}\\
\textit{1. Определим минимальное расстояние до точек} \\
\textit{2. ref - середина между предыдущей ref и точкой с минимальным расстоянием} \\
\textit{Таким образом мы способны отдалиться вдоль 2 осей на максимальное расстояние, не превышающее расстояние между точками в начале}


\begin{figure}[h]
	\centering
	\begin{tikzpicture}
		% Координаты вершин повёрнутого квадрата (45 градусов)
		\coordinate (A) at (0,1.414);
		\coordinate (B) at (1.414,0);
		\coordinate (C) at (0,-1.414);
		\coordinate (D) at (-1.414,0);
		\coordinate (P) at (0, 0);
		\coordinate (O) at (0, 2.814);
		\coordinate (N) at (0, 2.814/4);
		
		% Рисуем точки
		\fill[black] (A) circle (2pt);
		\fill[black] (B) circle (2pt);
		\fill[black] (C) circle (2pt);
		\fill[black] (D) circle (2pt);
		\fill[black] (P) circle (2pt);
		\fill[black] (O) circle (2pt);
		\fill[black] (N) circle (2pt);
		
		% Подписи
		\node[above right] at (A) {$point 1$};
		\node[below right] at (B) {$point 4$};
		\node[below left] at (C) {$point3$};
		\node[above left] at (D) {$point2$};
		\node[below] at (P) {$star, ref$};
		\node[above] at (O) {$G.C.$};
		\node[above right] at (N) {$ref'$};
	\end{tikzpicture}
	\caption{Примерное начальное состояние при $dist(star, point1) = d_{min}$. ref' - новая точка попадания ссылки после алгоритма определения ref}
	\label{fig:points}
\end{figure}
\subsection{Определение вектора минимума.}
Пусть есть некоторый вектор $\vec{a}$, который сонаправлен с вектором $\frac{df}{\vec{da}}$ градиента исследуемой функции f на минимум значения.

Определим градиент вдоль вектора $\vec{a}$ через его проекции:

\begin{equation}
	\Delta{a} = \frac{\Delta x_i} {\cos{\alpha_i}}
\end{equation}
где 
\begin{equation}
	\cos{\alpha_i} \propto \frac{\partial f}{\partial x_i}.
\end{equation}
Условие абсолютного минимума мы запишем как
\begin{equation}
	\frac{\Delta f}{\Delta a} = 0.
\end{equation}
Определим производную f по a как
\begin{equation}
	\frac{\Delta f}{\Delta a} = \sum_{i} \frac{\partial f}{\partial x_i} \cdot \frac{\partial x_i}{\partial a}
\end{equation}
\begin{equation}
	\Rightarrow \frac{\Delta f}{\Delta a} \propto \sum_{i} \frac{\partial f}{\partial x_i}^2,
\end{equation}
что значит
\begin{equation}
	\sum_{i} \frac{\partial f}{\partial x_i}^2 = 0.
\end{equation}
Поскольку каждая производная считается по производной по точке и при изменении остальных проекций, то
\begin{equation}
	\frac{\partial f}{\partial x_i} (\vec{x}) = \frac{\partial f}{\partial x_i} (\vec{x}_{0}) + \sum_{j} (\frac{\partial^2 f}{\partial x_i \partial x_j}(\vec{x}_0) \cdot \Delta x_j)
\end{equation}
Пренебрегая вторым порядком малости ($\Delta x_i \Delta x_j$), мы получаем упрощение:
\begin{equation}
	\sum_{i} \frac{\partial f}{\partial x_i}^2 = \sum_{i} J_i^2 + 2\sum_{i} (\Delta x_i \cdot \sum_{j} (J_j \cdot H_{ij}))
\end{equation}
Откуда мы можем записать это в более простом виде:
\begin{equation}
	A\cdot A^T + B\Delta{\vec{x}} = 0
\end{equation}
Применяя, что $B_{0i}= \sum_{i} H_{ij}\cdot J_j$, где J - матрица частных производных, а H - гессиан функций.
Из чего следует, что
\begin{equation}
	\Delta{\vec{x}} = -B^{-1}A 
\end{equation}

Также стоит обратить внимание на примитивный случай верного условия выполнения истинного минимума переменных:

$\bullet$ В такой верно подобранной точке (наборе аргументов $x_i$) будет наблюдаться, что
\begin{equation}
	\frac{df}{dx} = 0,
\end{equation}
из чего будет непременно следовать, что
\begin{equation}
	\frac{df}{dx_i}(\vec{x}_0) + \sum_{j}\frac{d^2 f}{dx_i dx_j}(\vec{x}_0)\Delta x_j = 0
\end{equation}
Решение этого уравнения лежит в однозначном определении переменных. Несомненно можно воспользоваться уравнением (1) (1-ое уравнение в данном подразделе) для какой либо догадки распределении величин $\Delta{x_i}/\Delta{x_j}$, но к результату оно приведёт неоднозначному, к тому же придётся решать квадратичное уравнение вида квадрата уравнения (7). Лёгкое упрощение можно заметить в том, где члены матрицы ковариаций дают наименьший вклад в текущую производную $\frac{df}{dx_i}$. Решение представит собой умножение обратной матрицы гессе с якобианом:
\begin{equation}
	\Delta{\vec{x}} = -H^{-1} J
\end{equation}
Исходя из решения действительной задачи на минимизацию коэффициента правдоподобия, решение уравнения (6) через библиотеки Python приводило к ухудшению значений функций и увеличению его в ~2 раза в зависимости от хода в вычислении производных (в любых случаях увеличивалось значение). Решение уравнения вида (13) давало сомнительный результат, который можно было бы едва заметить - менее 3 степени во множителе значения функции (почти нет). Таким образом вопрос стоит в двух вещах:

\textit{$\bullet$ Введении более точной функции минимизации, которая бы не была подвержена резким скачкам для выполнения условия уравнения (6)}

\textit{$\bullet$ Введение более точных методов определения минимума в функции множества переменных}\\
Последняя задача кажется невыполнимой в обычных случаях, но количество переменных всегда можно поменять в зависимости от требуемой точности модели (к примеру - 5 переменных разных численных порядков).

\subsection{Доверительные интервалы функции высоты поверхности $\zeta^n(X, Y)$.}
Python библиотекой scipy представляет устойчивое решение для значения $\mathcal{L}$
разными методами.
Исходя из этого, сама функция минимизации представляет собой устойчивое решение, что затрудняет определение у неё спонтанных отклонений в обе стороны. Исходя из этого, мы будем решать альтернативные решения для вектора $\vec{z}$ (альтернатива $\vec{x}$ ранее) путём ввода различных предварительных значений через минимизацию лограрифма модуля разностей коэффициентов правдоподобия в обоих случаях и выбирая максимальные отклонения по всем осям одновременно.
\begin{gather}
	\mathcal{L}(\vec{z}) \lesssim \mathcal{L}(\vec{z}_0)_0 \\
	|z_i - (z_{i})_0| \rightarrow max
\end{gather}
Таким образом мы получаем доверительные интервалы $\Delta \vec{z}$ для дальнейших манипуляций с ними.

Далее последует поиск в интервалах полученных значений $\Delta \vec{z}$ в проекциях для максимум отклонения функции $\zeta^n(X, Y)$ от полученных в первый раз значений на плоскости (X, Y):
\begin{equation}
	|\zeta^n(X, Y) - (\zeta^n(X, Y))_0| \rightarrow max
\end{equation}
\section*{Черновик}

Можно попробовать описать идею о том, что цефеиды имеют статистическое движение, характера уравнения несжимаемой жидкости, иначе эту идею можно представить в виде условия на разлёт звёзд:
\begin{equation}
	\sum_i \upsilon_{Ri} = 0.
\end{equation}
Т. е. в моделировании диска во времени должно соблюдаться это условие.


В простом случае мы можем распространить это выражение на долю наблюдаемых звёзд (связывая это с центральной симметрией линии Солнце-Центр Галактики) и, также, на текущий момент изменений расстояний, которое мы можем предсказать для ближайшего будущего, зная $\omega^2 R$ и $a_R$:
\begin{equation}
	v_R = v_{R_0} + (\frac{v_\tau^2}{R} - a_R)t + O(t^2),
\end{equation}
Пренебрегая вторым членом точности, получаем условие на весь диск или на некоторое расстояние $(R, R+\Delta R)$.

Возможно, есть смысл проверить на саму устойчивость данный подвижный диск, представляемый нами как звёздный, через определение точечных малых масс вдоль него и прослеживание за их радиальными движениями, считая их закреплёнными за дисками. Это проверит условие на возможность существования подобных искривлений в звёздном диске и поможет удостоверить нас в схожести реальной и наблюдаемой картины дисков.

Как это может быть реализовано: \\
\textit{$\bullet$ Мы разделим диск в каком-то месте на "части" от радиуса.} \\
\textit{$\bullet$ Будем смотреть на изменение потенциала вдоль изменения радиуса, а именно - поскольку $\vec{f} = -grad U$, то будем изучать направленность силы f, домножая на направление перпендекуляра к точке в плоскости} \\
Сам перпендекуляр будет иметь вид 
\begin{gather}
	\vec{\lambda} = \frac{da}{dx}\vec{e_x}+\frac{da}{dy}\vec{e_y}-\frac{da}{dz}\vec{e_z} \text{, где} \\
	a = -Z + \zeta^n(X, Y, z) \textit{, откуда} \\
	\frac{da}{dx} =  -\sum_i z_i kX^{k-1} Y^m \\
	\frac{da}{dx} =  -\sum_i z_i X^{k} mY^{m-1} \\
	\frac{da}{dz} = 1
\end{gather}
\end{document}